\begin{derivation}[从\eqnrefs{eq:sm-higgs-cartesian-rxi-dp11}到\eqnrefs{eq:sm-rxi-quadratic-before-gf-dp11}]
要得到 \(R_\xi\) 规范固定前必须消去的混合项，不能只读最终质量矩阵；需要把协变导数在真空附近按场次数展开。正文采用
\begin{equation*}
\Phi
=\sqrt{\hbar c}
\begin{pmatrix}
\widehat\phi^+\\[2pt]
\bigl(v+\widehat h+\ii\widehat\chi\bigr)/\sqrt2
\end{pmatrix}.
\end{equation*}
只保留 \(D_\mu\Phi\) 中对涨落场至多一次的部分。带电上分量可写成
\begin{equation*}
(D_\mu\Phi)_1^{(1)}
=\sqrt{\hbar c}\left(
\partial_\mu\widehat\phi^+
+\ii\kappa_W\mathcal W_\mu^+
\right),
\end{equation*}
其中 \(\kappa_W=m_Wc/\hbar\) 按正文定义。其模平方给
\begin{align*}
\frac{|(D_\mu\Phi)_1^{(1)}|^2}{\hbar c}
={}&\partial_\mu\widehat\phi^-\partial^\mu\widehat\phi^+
+\kappa_W^2\mathcal W^-_\mu\mathcal W^{+\mu}\\
&+\ii\kappa_W
\left(
\mathcal W^-_\mu\partial^\mu\widehat\phi^+
-\mathcal W^+_\mu\partial^\mu\widehat\phi^-
\right).
\end{align*}
因此带电 Goldstone 的动能、\(W\) 的质量二次型以及 \(W\)--Goldstone 导数混合来自同一个平方，三者并不是三个独立假设。

中性下分量的一阶部分同理为
\begin{equation*}
(D_\mu\Phi)_2^{(1)}
=\sqrt{\frac{\hbar c}{2}}
\left[
\partial_\mu\widehat h
+\ii\partial_\mu\widehat\chi
+\ii\kappa_Z\mathcal Z_\mu
\right],
\end{equation*}
其中真空方向对光子组合的系数为零，因此 \(A_\mu^{\rm geo}\) 不产生质量项。取模平方并只保留二次项，实部与虚部分离后得到
\begin{equation*}
\frac12(\partial\widehat h)^2
+\frac12(\partial\widehat\chi)^2
+\frac12\kappa_Z^2\mathcal Z_\mu\mathcal Z^\mu
+\kappa_Z\mathcal Z_\mu\partial^\mu\widehat\chi.
\end{equation*}
这一步同时说明为什么 \(Z\)--Goldstone 混合只涉及 \(\widehat\chi\)，而径向模 \(\widehat h\) 不与 \(Z_\mu\) 发生双线性导数混合。

最后处理势能。若真空由 \(V'(v)=0\) 固定，则在径向方向 Taylor 展开
\begin{equation*}
V(v+\widehat h)
=V(v)+\frac12V''(v)\widehat h^2
+\mathcal O(\widehat h^3),
\end{equation*}
一次项因驻点条件消失。用正文定义 \(\kappa_h^2\equiv V''(v)/(\hbar c)\) 后，拉氏量中的势能贡献为
\begin{equation*}
-\frac12\kappa_h^2\widehat h^2.
\end{equation*}
把带电、中性与径向三部分相加，就逐项得到正文\eqnrefs{eq:sm-rxi-quadratic-before-gf-dp11}。该二次核明确暴露出随后选择 \(R_\xi\) 规范固定函数所要消去的两类导数混合项。
\end{derivation}

\begin{derivation}[从\eqnrefs{eq:sm-rxi-gauge-functions-dp11}到\eqnrefs{eq:sm-rxi-goldstone-masses-dp11}]
先展开带电规范固定项的交叉部分：

\begin{align*}
 -\frac1{\xi_W}F^+F^-
 \supset{}&-\ii\kappa_W\widehat\phi^-\partial\!\cdot\mathcal W^+
 +\ii\kappa_W\widehat\phi^+\partial\!\cdot\mathcal W^-.
\end{align*}
对作用量积分分部，并假设场在积分边界足够快衰减，使表面项为零：
\begin{align*}
 \int\dd^4 x\,
 \widehat\phi^-\partial_\mu\mathcal W^{+\mu}
 &=-\int\dd^4 x\,
 (\partial_\mu\widehat\phi^-)\mathcal W^{+\mu},\\
 \int\dd^4 x\,
 \widehat\phi^+\partial_\mu\mathcal W^{-\mu}
 &=-\int\dd^4 x\,
 (\partial_\mu\widehat\phi^+)\mathcal W^{-\mu}.
\end{align*}
因此规范固定项贡献
\begin{equation*}
 +\ii\kappa_W\mathcal W^+\!\cdot\partial\widehat\phi^-
 -\ii\kappa_W\mathcal W^-\!\cdot\partial\widehat\phi^+,
\end{equation*}
恰与\eqnrefs{eq:sm-rxi-quadratic-before-gf-dp11} 的带电混合项逐项相消。中性部分为
\begin{equation*}
 -\frac1{2\xi_Z}(F^Z)^2
 \supset
 +\kappa_Z\widehat\chi\,\partial\!\cdot\mathcal Z
 \longrightarrow
 -\kappa_Z\mathcal Z\!\cdot\partial\widehat\chi,
\end{equation*}
同样消去 $Z$--$\chi$ 混合。

规范固定函数平方还留下纯标量项
\begin{equation*}
 -\xi_W\kappa_W^2\widehat\phi^+\widehat\phi^-
 -\frac12\xi_Z\kappa_Z^2\widehat\chi^2.
\end{equation*}
与标量传播子的标准二次质量项比较，得到
\begin{equation*}
 m_{\phi^\pm}^2=\xi_Wm_W^2,
 \qquad
 m_\chi^2=\xi_Zm_Z^2,
\end{equation*}
即\eqnrefs{eq:sm-rxi-goldstone-masses-dp11}。
\end{derivation}

\begin{derivation}[从\eqnrefs{eq:sm-rxi-gauge-functions-dp11}到\eqnrefs{eq:sm-rxi-ghost-masses-dp11}]
由正文的 $\mathcal M_{\rm FP}^{ab}(x,y)=\delta F^a(x)/\delta\alpha^b(y)$ 出发。在线性化真空附近，破缺规范变换使 Goldstone 场沿规范轨道移动：

\begin{equation*}
 \delta\widehat\phi^\pm
 =\pm\ii\kappa_W\alpha^\pm+\order(\text{场}),
 \qquad
 \delta\widehat\chi
 =\kappa_Z\alpha^Z+\order(\text{场}).
\end{equation*}
同时规范场的线性变换给
\begin{equation*}
 \delta(\partial\!\cdot\mathcal W^\pm)=\partial^2\alpha^\pm+\order(\text{场}),
 \qquad
 \delta(\partial\!\cdot\mathcal Z)=\partial^2\alpha^Z+\order(\text{场}),
\end{equation*}
以及 $\delta(\partial\!\cdot A^{\rm geo})=\partial^2\alpha^\gamma+\order(\text{场})$。代入\eqnrefs{eq:sm-rxi-gauge-functions-dp11}：
\begin{align*}
 \delta F^\pm
 &=\left(\partial^2+\xi_W\kappa_W^2\right)\alpha^\pm
 +\order(\text{场}),\\
 \delta F^Z
 &=\left(\partial^2+\xi_Z\kappa_Z^2\right)\alpha^Z
 +\order(\text{场}),\\
 \delta F^\gamma
 &=\partial^2\alpha^\gamma+\order(\text{场}).
\end{align*}
于是鬼场二次算符分别是
\begin{equation*}
 \partial^2+\xi_W\kappa_W^2,
 \qquad
 \partial^2+\xi_Z\kappa_Z^2,
 \qquad
 \partial^2,
\end{equation*}
与自由标量二次核比较便得到\eqnrefs{eq:sm-rxi-ghost-masses-dp11}。
\end{derivation}

\begin{derivation}[从\eqnrefs{eq:sm-rxi-gf-lagrangian-dp11}到\eqnrefs{eq:sm-rxi-massive-vector-propagator-dp41}]
取物理四动量 $p^\mu$，并记

\begin{equation*}
 P^T_{\mu\nu}=\eta_{\mu\nu}-\frac{p_\mu p_\nu}{p^2},
 \qquad
 P^L_{\mu\nu}=\frac{p_\mu p_\nu}{p^2}.
\end{equation*}
它们满足
\begin{equation*}
 (P^T)^2=P^T,
 \qquad
 (P^L)^2=P^L,
 \qquad
 P^TP^L=0,
 \qquad
 P^T+P^L=I.
\end{equation*}
规范固定后的约化二次核可写成
\begin{equation*}
 K_{\mu\nu}(p)
 =-(p^2-m_V^2c^2)\eta_{\mu\nu}
 +\left(1-\frac1{\xi_V}\right)p_\mu p_\nu.
\end{equation*}
用 $\eta=P^T+P^L$ 与 $p_\mu p_\nu=p^2P^L_{\mu\nu}$ 分解：
\begin{align*}
 K_{\mu\nu}
 ={}&-(p^2-m_V^2c^2)P^T_{\mu\nu}\\
 &-\frac1{\xi_V}(p^2-\xi_Vm_V^2c^2)P^L_{\mu\nu}.
\end{align*}
两个投影子空间互相正交，所以逆核逐块取倒数：
\begin{equation*}
 K^{-1}_{\mu\nu}
 =-\frac{P^T_{\mu\nu}}{p^2-m_V^2c^2}
 -\frac{\xi_VP^L_{\mu\nu}}{p^2-\xi_Vm_V^2c^2}.
\end{equation*}
乘上 Feynman 边界条件对应的 $\ii$ 并重新合并两个投影算符，得到
\begin{equation*}
 D^{(V)}_{\mu\nu}(p)
 =\frac{-\ii}{p^2-m_V^2c^2+\ii0^+}
 \left[
 \eta_{\mu\nu}
 -(1-\xi_V)
 \frac{p_\mu p_\nu}{p^2-\xi_Vm_V^2c^2+\ii0^+}
 \right],
\end{equation*}
即\eqnrefs{eq:sm-rxi-massive-vector-propagator-dp41}。$\xi_V=1$ 时纵向附加项消失；$\xi_V\to\infty$ 时传播子化为幺正规范形式。这两个极限分别检验了纵向极点结构与质量项的归一化。
\end{derivation}
